\chapter{结论与展望}
系统在Nexys A7-100T上实现了声音、加速度、温度、片内电压和GPIO采集，主机软件集成实时监测、事件记录、听歌识曲与体感游戏。FPGA负责接口时序、特征计算和缓存，主机负责文件处理、检索与交互，各应用共用一个串口服务。

双组事件缓存保存触发前0.5秒和触发后1秒的数据，独立音乐缓冲支持20秒、8 kHz录音。31抽头FIR及二倍抽取将长录音容量控制在160000点；事件和音乐均通过编号读取、整段校验及确认释放管理存储。三敲触发、任务中止和模式恢复支持识曲、监测与游戏之间的切换。

布线后电路满足100 MHz约束，BRAM使用119.5/135，DSP使用9/240。系统完成完整事件导出、三首录音的在线识曲回放和八关实板游戏测试，独立EXE能够在中文路径及独立数据目录下启动并完成资源校验。

后续工作包括传感器标定、跨设备测试及扩大独立测试样本。识别算法可增加声音事件分类和学习式声振融合，并比较检出率、误触发、资源占用及延迟；硬件可扩展外部存储和更高速通信，以增加录音长度并缩短导出时间。

